% concatenated from pir.tex

\documentclass[12pt,a4paper]{amsart}
\usepackage{amsaddr}
\usepackage{amsmath,amscd}
\usepackage{latexsym,amsmath,amssymb,mathrsfs,setspace,enumerate,tikz}
\usepackage{graphicx}
\usepackage{subcaption}
\usepackage[T1]{fontenc}
\usepackage[utf8]{inputenc}
\usepackage{lmodern}
\usepackage{amssymb}
\usepackage{multicol}
\usepackage{tikz-network}
\usetikzlibrary{automata,positioning}
\usepackage[english]{babel} 
\usepackage{blindtext}

\theoremstyle{plain}
\newtheorem{theorem}{Theorem}[section]
\newtheorem{lemma}[theorem]{Lemma}
\newtheorem{proposition}[theorem]{Proposition}
\newtheorem{corollary}[theorem]{Corollary}
\theoremstyle{definition}
\newtheorem{definition}[theorem]{Definition}
\newtheorem{notation}[theorem]{Notation}
\newtheorem{remark}[theorem]{Remark}
\newtheorem{example}[theorem]{Example}
\newtheorem*{note}{Note} 
\DeclareMathOperator{\dom}{dom}
\DeclareMathOperator{\fix}{fix}

\newcommand{\N}{\mathbb{N}}
\newcommand{\B}{\mathbf{B}}
\newcommand{\inv}[2]{#1^{-1}(#2)}
\begin{document}
	
	\title[Ordered Semigroups and categories of ideals]{Ordered semigroups and ideals categories of  principal ideal rings}
	
			\author[ P. K. Minnumol and P. G. Romeo]{P. K. Minnumol and P. G. Romeo}
			
	\address{Department of Mathematics\\
		Cochin University of Science and Technology(CUSAT)\\Kochi, Kerala, India, 682022}
	\email{$pkminnumol@gmail.com,\, romeo_-parackal@yahoo.com$}
	
%\author[Minnumol P K]{Minnumol P K}
%\address{Research Scholar\\Department of Mathematics\\
%	Cochin University of Science and Technology(CUSAT)\\Kochi, Kerala, India, 682022}
	\thanks{First author wishes to thank Cochin University of Science And Technology for providing financial support.\\
		Second author wishes to thank Council of Scientific and Industrial Research(CSIR) INDIA, for providing financial support.
		 }
	\subjclass{13F10, 20M17, 20M14, 06F05}
	\keywords{ordered semigroup, ideals, principal ideal ring, category}

\begin{abstract}
	In this paper, we view the collection of ideals of a commutative principal ideal ring from two perspectives: one as an ordered semigroup $\mathcal{I}(R)$ and the other as a category $\mathbb{I}_R$. It is shown that $\mathcal{I}(R)$ is a regular ordered semigroup whereas $\mathbb{I}_R$ is a category with subobjects. Further we establish the correspondence between these structures.
\end{abstract}
\maketitle
An ordered semigroup is a semigroup equipped with a compatible partial order, there are various investigations carried out in ordered semigroups. 
Some authors have studied partial orders on regular semigroups; an ordered semigroup of this type is called an \emph{ordered regular semigroup}. 
On the other hand, some authors have studied regularity within ordered semigroups; such a notion is referred to as \emph{ordered regularity}, and semigroups in which every element is ordered regular are called \emph{regular ordered semigroups}. 
Since in an ordered semigroup every regular element is ordered regular the class of regular ordered semigroups is much larger than the class of ordered regular semigroups. 
Here we construct the ordered semigroup structure as well as the categorical structure of ideals of a commutative principal ideal ring.  We establish the correspondence between these structures.

\section{Preliminaries}
First we recall some definitions and basic results regarding ordered semigroups, rings and category theory needed in this paper.
\subsection{Ordered semigroups}

An \textit{ordered semigroup}  $(S,\cdot,\leq)$ is a partially ordered set $(S,\leq)$ together with a semigroup structure $(S,\cdot)$ such that the order is compatible with multiplication; \cite{txtbk}. 
For a subset $A \subseteq S$, the set
\[
(A] = \{x \in S \mid (\exists a \in A)\; x \leq a\}
\]
is called the \textit{downward closure} of $A$. A subset $A$ is said to be \textit{downward closed} if $(A] = A$.

\begin{definition}\cite{keharegularity}
	An element $a \in S$ is said to be \textit{ordered regular} if there exists $x \in S$ such that
	$a \leq axa$ or equivalently $a \in (aSa]$, ordered regular elements of $S$ is denoted by $Reg_{\leq}(S)$
\end{definition}
When $Reg_{\leq}(S)=S$, then $S$ is called a \textit{regular ordered semigroup}.
\begin{example}
Consider $(\mathbb{N},\cdot,\leq)$, where $\mathbb{N}$ denotes the set of natural numbers with the usual multiplication and order. Clearly, $1$ is the only regular element in the classical sense, but for every $n \in \mathbb{N}$, there exists $x \in \mathbb{N}$ such that $n \leq nxn$, hence, every element of $\mathbb{N}$ is ordered regular and $(\mathbb{N},\cdot,\leq)$ is a regular ordered semigroup.
\end{example}

An ordered semigroup $S$ is called \textit{intra-regular}  if for each $a \in S$ there exist $x,y \in S$ such that
$a \leq xa^{2}y$
or equivalently $a \in (Sa^{2}S]$, \cite{intrregular}.
An element $a' \in S$ is said to be an \textit{ordered inverse}  of $a$ if
$a \leq aa'a \quad \text{and} \quad a' \leq a'aa'.$
The set of all ordered inverses of $a$ is denoted by $V_{\leq}(a)$, \cite{jamadar}.
An ordered semigroup $S$ is called an \textit{inverse ordered semigroup} if for each $a \in S$, any two ordered inverses of $a$ are $\mathcal{H}$-related.
An element $e \in S$ is called an \textit{ordered idempotent} if $e \leq e^{2}$.

The following are some important classes of regular ordered semigroups.
\begin{itemize}
	\item an ordered semigroup $S$ is \textit{group-like} if for all $a,b \in S$ there exist $x,y \in S$ such that 	$a \leq xb$ and $a \leq by$.
	\item an ordered semigroup $S$ is \textit{completely regular} if for every $a \in S
	a \in (a^{2}Sa^{2}].$
	\item a regular ordered semigroup $S$ is called a \textit{Clifford ordered semigroup} if for all $a \in S$ and $e \in E_{\leq}(S)$ there exist $u,v \in S$ such that
	$ae \leq eua \quad \text{and} \quad ea \leq ave.$
	
\end{itemize} 
A nonempty subset $A$ of $S$ is called a \textit{left ideal} ( \textit{right ideal}) of $S$ if
$SA \subseteq A$ ($AS \subseteq A$) and $(A]=A$.
A subset $A$ is called an \textit{ideal} of $S$ if it is both a left and a right ideal.
The principal left ideal, principal right ideal and principal ideal generated by $a \in S$ are denoted by $L(a)$, $R(a)$, and $I(a)$, respectively and are defined in the usual manner \cite{kahaloideal}.
\begin{center}
	$ L(a) = (a \cup Sa] $\\
	$ R(a) = (a \cup aS] $\\
	$ I(a) = (a \cup Sa \cup aS \cup SaS] $\\
	
\end{center}
The \textit{Green's relations}  $ \mathscr{L} ,\mathscr{R}, \mathscr{J}, \mathscr{H}, \mathscr{D}  $ on an ordered semigroup $ S $ are given as follows:
\begin{center}
	$ a\,\mathscr{L} \,b $ if and only if $ L(a) = L(b) $\\
	$ a\,\mathscr{R} \,b $ if and only if $ R(a) = R(b) $\\
	$ a\,\mathscr{J} \,b $ if and only if $ I(a) = I(b) $\\
	$ \mathscr{H} = \mathscr{L} \cap \mathscr{R} $\\
	$ \mathscr{D} = \mathscr{L} \vee \mathscr{R} $
\end{center}
If $ S $ is a regular ordered  semigroup for each $ a \in{S} $ there exists $ x \in{S} $ such that $ a \leq axa$, hence $ a \in (Sa] $ and 
$ a \in (aS] $ and so the Green's equivalence for regular ordered semigroups are:
\begin{center}
	$ a\,\mathscr{L} \,b $ if and only if $ (Sa]  = (Sb] $\\
	$ a\,\mathscr{R} \,b $ if and only if $ (aS]  = (bS] $\\
\end{center}
The following theorem characterizes Green’s relations in a regular ordered semigroup.
\begin{proposition}(cf.\cite{mpk})
	Let $ a $ and $ b $ be elements of a regular ordered semigroup $ S $. Then 	$ a\,\mathscr{L} \,b $ if and only if there exist $ x,y \in{S} $ such that $ a \leq xb $ and $ b \leq ya $. Also $ a\,\mathscr{R} \,b $ if and only if there exist $ u,v \in{S} $ such that $ a \leq bu $ and $ b \leq av $.
\end{proposition}
The inverse transversals for regular ordered semigroups. is defined below.  
\begin{definition}\label{definvtrsl}(cf.\cite{mpkinverse})
	Let $(S,\cdot,\leq)$ be a regular ordered semigroup. A subset $S_{0}\subseteq S$
	is called an inverse transversal\index{Inverse transversal} of $S$ if the following conditions hold:
	\begin{enumerate}
		\item $S_{0}$ is an inverse ordered subsemigroup of $S$;
		\item $(S_{0}] = S_{0}$;
		\item $V_{\le}(a)\cap S_{0} \neq \varnothing$ for every $a \in S$;
		\item if $a',a'' \in V_{\le}(a)\cap S_{0}$, then $a' \,\mathcal{H}\, a''$ in $S_{0}$.
	\end{enumerate}
\end{definition}
\subsection{Rings and ideals}
A ring is a set $R$ equipped with two operations  addition $+$ and  multiplication $\cdot$ such that $(R,+)$ is an abelian group and $(R,\cdot)$ is a semigroup and multiplication distributes over addition. 
A subset $I$ of $R$ is called a \emph{left (right) ideal} of $R$ if $I$ is a subgroup of $(R,+)$ and $RI \subseteq I (IR \subseteq I)$ and a \emph{two-sided ideal} if it is both a left
ideal and a right ideal of $R$ and an ideal (left, right, or two-sided) $I$ of a ring $R$ is said to be finitely generated if there exists a finite subset $X=\{x_1,\dots,x_n\}$ of $I$ such that $I$ is generated by $X$. When $I$ is generated by a single element then it is called a principal ideal.
\begin{definition}
	A ring $R$ is called Noetherian if every ideal of $R$ is finitely generated.
\end{definition}
\begin{definition}(cf.\cite{musili})
	A ring $R$ is called a principal ideal ring (PIR)\index{Principal ideal ring} if every ideal of $R$ is principal.
\end{definition}
\subsection{Category theory}	
	A \textit{category} $ \mathcal C $ consisting of a class of  objects  written as $\nu {\mathcal C}$  and collection of morphisms $ f \in \mathcal C(A,B)$ from each object $ A = dom\; f $  to each object $ B = cod\; f $. For each pair $ (f,g) $ of morphisms with $ dom\; g=cod\; f $, a morphism $fg : \dom f \, \rightarrow\,cod\; g$ is the composition  and for each object $ A $ there exist a  unique morphism   $  1_A \in \mathcal C( A,A ) $ is  called the identity morphism on $A$. Further the composition satisfies $ f(gh) = (fg)h $ whenever defined and $  1_A f = f = f 1_B  $ for all $f \in  \mathcal C ( A,B) $.\\
The maps between categories are called functors. For categories $\mathcal{C}$ and $ \mathcal{D}  $
	a \emph{contravariant functor} $F : \mathcal{C} \to \mathcal{D}$ assigns to each object 
$A \in \nu\mathcal{C}$ an object $F(A) \in \nu\mathcal{D}$ and to every morphism 
$f : A \to B$ in $\mathcal{C}$ a morphism
$F(f) : F(B) \to F(A)$
in $\mathcal{D}$.
These assignments satisfy the following conditions:
\begin{itemize}
	\item For every object $A$ of $\mathcal{C}$,
	$	F(1_A) = 1_{F(A)}.$
	
	\item For any pair of composable morphisms $f : A \to B$ and $g : B \to C$ in $\mathcal{C}$,
	\[
	F(fg) = F(g) F(f) .
	\]
\end{itemize}
\begin{definition}
	A morphism \(f : A \to B\) is called a \emph{monomorphism}\index{Monomorphism} if, for any
	morphisms \(h',k' : D \to A\), the condition \(h'f = k'f\) implies \(h' = k'\).
	Thus monomorphisms are precisely the right-cancellable morphisms.
\end{definition}
A preorder $\mathcal {P}$ is a category such that for any $p, p' \in \nu \mathcal {P}$, the
hom-set $\mathcal {P}(p,p')$ contains at most one morphism. In this case, there is a
quasi-order relation $\preceq$ on $\nu \mathcal {P}$ defined by
\[
p \preceq p' \;\Longleftrightarrow\; \mathcal {P}(p,p') \neq \varnothing.
\]
The preorder $P$ is said to be \emph{strict} if $\preceq$ is a partial order (see cf.~\cite{kss}).
\begin{definition}(cf.\cite{kss})\label{suboject}
	Let $\mathcal{C}$ be a category and $\mathcal{P}$ a subcategory of $\mathcal{C}$.
	The pair $(\mathcal{C}, \mathcal{P})$ is called a category with subobjects\index{Subobject category}
	if the following conditions hold:
	\begin{itemize}
		\item $\mathcal{P}$ is a strict preorder with $\nu\mathcal{C} = \nu\mathcal{P}$;
		\item every morphism in $\mathcal{P}$ is a monomorphism;
		\item if $f, g \in \mathcal{P}$ and $f =  hg $ for some $h \in \mathcal{C}$,
		then $h \in \mathcal{P}$.
	\end{itemize}
	For objects $C, D \in \nu\mathcal{C}$, we denote the unique morphism in
	$\mathcal{P}$ from $C \to D$ by $j_{(C, D)}$, called the inclusion.
	In this case $C$ is referred to as a subobject of $D$.
	
\end{definition}
%\begin{definition}(cf.\cite{kss})	Let $\mathcal{C}$ be a category with subobjects.	A canonical factorization of a morphism $f$ in $\mathcal{C}$	is a factorization of the form	\[	f = qj ,	\]	where $q$ is an epimorphism and $j$ is an inclusion.\end{definition}	

\section{Ordered semigroup of ideals of a principal ideal ring.}\label{section}

In this section, we establish that the set of all ideals of a commutative principal ideal ring carries a natural structure of a regular ordered semigroup, where the semigroup operation is ideal multiplication and the partial order is determined by the divisibility of generators. This order is compatible with the multiplication and every element of the semigroup is ordered regular.

Let $R$ be a commutative principal ideal ring with unity and $\mathcal{I}(R)$ the set of all ideals of $R$. 
Then for every $A \in \mathcal{I}(R)$, there exists $a \in R$ such that
\[
A = \langle a \rangle = \{ra \mid r \in R\}.
\]
For $A, B \in \mathcal{I}(R)$, the product of ideals $A$ and $B$ is defined by
\[
AB = \left\{ \sum_{i=1}^{n} a_i b_i \;\middle|\; a_i \in A,\ b_i \in B,\ n \in \mathbb{N} \right\},
\]
clearly, equipped with this operation $\mathcal{I}(R)$ is a semigroup.
For, consider principal ideals $A = \langle a \rangle$ and $B = \langle b \rangle$, then
\[
\langle a \rangle \langle b \rangle 
= \left\{ \sum_{i=1}^{n} r_i a\, s_i b \;\middle|\; r_i, s_i \in R,\ n \in \mathbb{N} \right\}.
\]
When $R$ is commutative the product of principal ideals simplifies to
\[
AB = \langle a \rangle \cdot \langle b \rangle = \langle ab \rangle.
\]
Indeed, for any $r \in R$, we have $r(ab) = (ra)b \in \langle a\rangle\langle b\rangle$, which shows that
$\langle ab\rangle \subseteq \langle a\rangle\langle b\rangle$.
On the other hand, each element
$ x \in \langle a\rangle\langle b\rangle $
can be written as a finite sum $\sum_{i=1}^n (r_i a)(s_i b)$ with $r_i,s_i \in R$,
 hence $x \in \langle ab\rangle$, and,
$\langle a\rangle\langle b\rangle = \langle ab\rangle.$

For $A  = \langle a \rangle, B = \langle b \rangle \in \mathcal{I}(R)$, define the order relation $\preccurlyeq$ as follows: 
\[
A \preccurlyeq B \iff a \mid b \;\text{in}\; R,
\]
then $\preccurlyeq$ is a partial order on $\mathcal{I}(R)$, and is compatible with multiplication. For,  $A = \langle a \rangle$, since $a \mid a, \quad a = a \cdot 1$ and so $A \preccurlyeq A$.
Suppose $A \preccurlyeq B$ and $B \preccurlyeq A$, then $a \mid b$ and $b \mid a$, ie., there exist
$x,y \in R$ such that $b=ax$ and $a=by$. It follows that
$\langle a\rangle \subseteq \langle b\rangle$ and
$\langle b\rangle \subseteq \langle a\rangle$ and hence
$\langle a\rangle=\langle b\rangle$, and so $A =B$.
When $A \preccurlyeq B$ and $B \preccurlyeq C$, where $C = \langle c \rangle$, then $a \mid b$ and $b \mid c$, so there exist $x, y \in R$ such that
$b = ax $ and $ c = by,$  hence $a \mid c$ and therefore  $A \preccurlyeq C$, ie., the relation 
$\preccurlyeq$ is a partial order on $\mathcal{I}(R)$.
Also since $R$ is commutative, it easily follows that $\preccurlyeq$ is compatible with the semigroup multiplication on $\mathcal{I}(R)$.

Finally, we establish ordered regularity. Let $A \in \mathcal{I}(R)$. Then there exists $a \in R$ such that $A = \langle a \rangle$. 
For any $x \in R$, we have $a \mid axa$. Let $X = \langle x \rangle$. Then $
\langle axa \rangle = \langle a \rangle \langle x \rangle \langle a \rangle = AXA.$
Thus $A \preccurlyeq AXA$.
Therefore, each element of $\mathcal{I}(R)$ is ordered regular. Hence the triple $(\mathcal{I}(R), \cdot, \preccurlyeq)$ forms a regular ordered semigroup.
The preceding results are collected in the following theorem.

\begin{theorem}
	Let $R$ be a commutative principal ideal ring with unity and $\mathcal{I}(R)$ denote the set of all ideals of $R$. 
	Define a binary operation on $\mathcal{I}(R)$ by ideal multiplication and define a relation $\preccurlyeq$ on $\mathcal{I}(R)$ by for  $A  = \langle a \rangle, B = \langle b \rangle \in \mathcal{I}(R)$
	\[
	A \preccurlyeq B \iff a \mid b \; \text{in}\; R,
	\]
	Then $\big(\mathcal{I}(R), \cdot, \preccurlyeq\big)$ is a regular ordered semigroup.
\end{theorem}
\begin{remark}
	Suppose that $A, B \in \mathcal{I}(R)$ admit multiple generators, say 
	$A = \langle a \rangle = \langle a' \rangle$ and 
	$B = \langle b \rangle = \langle b' \rangle$. 
	If $a \mid b$ then $a' \mid b'$.  Suppose that $a \mid b$, then $b = ra$ for some $r \in R$, since $b' \in \langle b \rangle$, there exists $x \in R$ such that $b' = xb$. 
	Also, $a \in \langle a' \rangle$, implies there exists $y \in R$ such that $a = ya'$ and hence,
	\[
	b' = xb = x(ra) = xr(ya') = (xry)a',
	\]
	which shows that $a' \mid b'$.
\end{remark}
\begin{remark} When the ring $R$ is not commutative 
$\langle a\rangle \langle b\rangle = \langle ab\rangle$
does not  hold. 
\end{remark}
\begin{example}
	Consider the ring 
	\[
	R=\left\{ 
	\begin{pmatrix}
		a & b\\
		0 & c
	\end{pmatrix}
	\;\middle|\; a,b,c \in F
	\right\}
	\]
	of all upper triangular $2\times 2$ matrices over a field $F$ and the two sided ideals
	\[
	I_{1}
	= \left\langle 
	\begin{pmatrix}
		1 & 0\\
		0 & 0
	\end{pmatrix}
	\right\rangle
	= \left\{
	\begin{pmatrix}
		a & b\\
		0 & 0
	\end{pmatrix}
	\;\middle|\; a,b \in F
	\right\}
	\]
	and
	\[
	I_{2}
	= \left\langle 
	\begin{pmatrix}
		0 & 0\\
		0 & 1
	\end{pmatrix}
	\right\rangle
	= \left\{
	\begin{pmatrix}
		0 & c\\
		0 & d
	\end{pmatrix}
	\;\middle|\; c,d \in F
	\right\}, 
	\]
	
	then
	\[
	I_{1} I_{2}
	= \left\langle 
	\begin{pmatrix}
		1 & 0\\
		0 & 0
	\end{pmatrix}
	\right\rangle
	\left\langle 
	\begin{pmatrix}
		0 & 0\\
		0 & 1
	\end{pmatrix}
	\right\rangle
	= \left\{
	\begin{pmatrix}
		0 & a\\
		0 & 0
	\end{pmatrix}
	\;\middle|\; a \in F
	\right\}
	= \left\langle
	\begin{pmatrix}
		0 & 1\\
		0 & 0
	\end{pmatrix}
	\right\rangle.
	\]
	
	On the other hand,
	\[
	\left\langle
	\begin{pmatrix}
		1 & 0\\
		0 & 0
	\end{pmatrix}
	\begin{pmatrix}
		0 & 0\\
		0 & 1
	\end{pmatrix}
	\right\rangle
	=
	\left\langle
	\begin{pmatrix}
		0 & 0\\
		0 & 0
	\end{pmatrix}
	\right\rangle,
	\]
	
	hence,
	\[
	\left\langle
	\begin{pmatrix}
		1 & 0\\
		0 & 0
	\end{pmatrix}
	\right\rangle
	\left\langle
	\begin{pmatrix}
		0 & 0\\
		0 & 1
	\end{pmatrix}
	\right\rangle
	\neq
	\left\langle
	\begin{pmatrix}
		1 & 0\\
		0 & 0
	\end{pmatrix}
	\begin{pmatrix}
		0 & 0\\
		0 & 1
	\end{pmatrix}
	\right\rangle.
	\]
\end{example}

Next we investigate some properties of regular ordered semigroups of ideals of commutative principal ideal rings.

\begin{theorem}
	Let $R$ and $S$ be commutative principal ideal rings with unity.
	If $R$ and $S$ are isomorphic then the semigroups
	$(\mathcal{I}(R), \cdot , \preccurlyeq_{R})$ and
	$(\mathcal{I}(S), \cdot , \preccurlyeq_{S})$ are isomorphic.
\end{theorem}

\begin{proof}
	Let $\phi \colon R \to S$ be a ring isomorphism. Define a map 
	$\psi \colon \mathcal{I}(R) \to \mathcal{I}(S)$ by
	\[
	\psi(A) = \langle \phi(a) \rangle , \quad \text{for } A = \langle a \rangle \in \mathcal{I}(R).
	\]
	To show that $\langle \phi(a) \rangle = \phi(\langle a \rangle)$, consider 
	$y \in \langle \phi(a) \rangle$  then $y = s\,\phi(a)$ for some $s \in S$. 
	Since $\phi$ is surjective, there exists $r \in R$ such that $s = \phi(r)$ and so 
	$y = \phi(r)\phi(a) = \phi(ra).$ Also since $ra \in \langle a \rangle$, we have $y \in \phi(\langle a \rangle)$  thus $\langle \phi(a) \rangle \subseteq \phi(\langle a \rangle)$.
	
	Conversely, for $x \in \phi(\langle a \rangle),\quad x = \phi(ra)$ for some $r \in R$, and so $	x = \phi(r)\phi(a),$
	which implies $x \in \langle \phi(a) \rangle$. Thus 
	$\phi(\langle a \rangle) \subseteq \langle \phi(a) \rangle$, and therefore
$	\langle \phi(a) \rangle = \phi(\langle a \rangle).$
	
	Next, suppose $\psi(A) = \psi(B)$, then
$	\langle \phi(a) \rangle = \langle \phi(b) \rangle$
	implies $\phi(\langle a \rangle) = \phi(\langle b \rangle)$. 
	Since $\phi$ is injective it follows that $\langle a \rangle = \langle b \rangle$, 
	so $A = B$, ie., $\psi$ is injective. For $T = \langle t \rangle \in \mathcal{I}(S)$, 
	since $\phi$ is surjective, there exists $d \in R$ such that $t = \phi(d)$
	showing that $\psi$ is surjective. and hence $\psi$ is bijective.
	
	To verify that $\psi$ is a semigroup homomorphism, consider $A = \langle a \rangle,\, B = \langle b \rangle\, \in \mathcal{I}(R)$; 
	\[
	\psi(AB)
	= \psi(\langle a \rangle \langle b \rangle)
	= \psi(\langle ab \rangle)
	= \langle \phi(ab) \rangle
	= \langle \phi(a)\phi(b) \rangle
	= \psi(A)\psi(B),
	\]
	and so $\phi(ab) = \phi(a)\phi(b)$.
	Finally, $\psi$ preserves the order, for $A = \langle a \rangle\, B = \langle b \rangle,\, \in \mathcal{I}(R)$ such that 
	$A \preccurlyeq_R B$. Then $a \mid b$, so $b = ax$ for some $x \in R$, hence $\phi(b) = \phi(ax) = \phi(a)\phi(x),$
	which shows that $\phi(a) \mid \phi(b)$, ie., $	\langle \phi(a) \rangle \preccurlyeq_S \langle \phi(b) \rangle,$ implies, $\psi(A) \preccurlyeq_S \psi(B).$
	
	Therefore, $\psi$ is a bijective order-preserving semigroup homomorphism, consequently,
	\[
	(\mathcal{I}(R), \cdot, \preccurlyeq_R)
	\cong
	(\mathcal{I}(S), \cdot, \preccurlyeq_S).
	\]
\end{proof}

\begin{remark}
The reverse implication of the preceding theorem need not hold in general.
To illustrate, consider the fields $\mathbb{R}$ and $\mathbb{Q}$, each one is a commutative principal ideal ring with identity whose only ideals of are
$\langle 0 \rangle $ and $\langle 1 \rangle$.
As a consequence the associated ordered semigroups of ideals are isomorphic
\[
(\mathcal{I}(\mathbb{R}), \cdot, \preccurlyeq_{\mathbb{R}})
\cong
(\mathcal{I}(\mathbb{Q}), \cdot, \preccurlyeq_{\mathbb{Q}})
\]
nevertheless the rings $\mathbb{R}$ and $\mathbb{Q}$ themselves are not isomorphic.
\end{remark}
An inverse transversal of the ordered semigroup of ideals of $R$ is identified below.

\begin{proposition}
	Let $R$ be a commutative principal ideal ring with unity and 
	$\big(\mathcal{I}(R), \cdot, \preccurlyeq\big)$ be the ordered semigroup of
	ideals. Then $\{ \langle 1\rangle \}$ is an inverse transversal of
	$\big(\mathcal{I}(R), \cdot, \preccurlyeq\big)$.
\end{proposition}

\begin{proof}
	Consider $ A = \langle a\rangle \in \mathcal{I}(R)$. Since $R = \langle 1 \rangle$ and 
	$a = a \cdot 1 \cdot a$, it follows that $a \mid a1a$, hence $A \preccurlyeq ARA.$
	Similarly, $1 \mid 1a1$, which implies $R \preccurlyeq RAR,$ thus $R$ is as an ordered inverse for every element of $\mathcal{I}(R)$. 	
	Moreover, $	RR = R \quad \text{and} \quad R \preccurlyeq R,$
	so $R$ forms an inverse ordered subsemigroup of $\big(\mathcal{I}(R), \cdot, \preccurlyeq\big)$. Clearly $(R] = \{R\},$ and all the conditions for an inverse transversal (see Definition~\ref{definvtrsl}) are satisfied. Therefore,
$	\mathcal{I}(R)_0 = \{R\}.$
\end{proof}

We next examine the connection between regularity of commutative
principal ideal ring and regularity of its semigroup of ideals.

\begin{proposition}
	Let $R$ be a commutative principal ideal ring with unity and $a \in R$.
	Then $a$ is von Neumann regular in $R$ if and only if
	$\langle a \rangle$ is a regular\index{Regular} element of the semigroup
	$\mathcal{I}(R)$.
\end{proposition}
\begin{proof}
	Let $a$ be a von Neumann regular element of $R$. Then there exists $x \in R$ such that
$	a = axa.$
	Consequently $	\langle a\rangle = \langle axa\rangle
	= \langle a\rangle \langle x\rangle \langle a\rangle$
	which shows $\langle a\rangle$ is a regular element of the semigroup
	$\mathcal{I}(R)$.\\
		
	\noindent
	Conversely suppose that $\langle a\rangle$ is regular in $\mathcal{I}(R)$.
	Then there exists $x \in R$ satisfying
$	\langle a\rangle
	= \langle a\rangle \langle x\rangle \langle a\rangle
	= \langle axa\rangle.$
	Since $R$ is a commutative principal ideal ring, equality of principal ideals yields
$	a = r(axa) = a(rx)a
	\quad \text{for some } r \in R.$
Hence $a$ is von Neumann regular in $R$.
\end{proof}
In the following theorem we investigate Green’s relations\index{Green’s relations}
on the ordered semigroup $\big(\mathcal{I}(R), \cdot , \preccurlyeq\big)$.
\begin{theorem}
	Let $R$ be a commutative principal ideal ring with unity and 
	$\big(\mathcal{I}(R), \cdot , \preccurlyeq\big)$  the ordered semigroup of ideals of $R$.
	Then Green’s relations $\mathcal{L}, \mathcal{R}, \mathcal{H}, \mathcal{D}$, and $\mathcal{J}$ 	on $\mathcal{I}(R)$ are all universal relations.
\end{theorem}

\begin{proof}
	The existence of the zero ideal together with the
	commutativity of $R$ forces all Green’s relations on
	$\mathcal{I}(R)$ to be universal. For, consider the ordered semigroup
	$\big(\mathcal{I}(R), \cdot , \preccurlyeq\big)$, since $a \mid 0$ for every $a \in R$, it follows that
$\langle a \rangle \preccurlyeq \langle 0 \rangle \quad \text{for all } a \in R.$
Thus for every $A \in \mathcal{I}(R)$, we have $A \preccurlyeq 0$, where $0 = \langle 0 \rangle$.
	
Let $(\mathcal{I}(R)A]$ denote the ideal of the ordered semigroup $\mathcal{I}(R)$ generated by $A$. Since $0 \in \mathcal{I}(R)$ and
$0A = 0,$ 	it follows that $0 \in (\mathcal{I}(R)A]$. 
	As ideals in an ordered semigroup are downward closed with respect to $\preccurlyeq$ and since $A \preccurlyeq 0$ for all $A \in \mathcal{I}(R)$, we obtain
$(\mathcal{I}(R)0] = \mathcal{I}(R).$
	Hence, the ordered semigroup $\big(\mathcal{I}(R), \cdot , \preccurlyeq\big)$ admits no proper ideals.
	Since $R$ is commutative the semigroup $\big(\mathcal{I}(R), \cdot\big)$ is also commutative, therefore, all Green’s relations coincide; that is,
	\[
	\mathcal{L} = \mathcal{R} = \mathcal{J} = \mathcal{H} = \mathcal{D}
	\quad \text{on } \mathcal{I}(R).
	\]
	Moreover, for any $A \in \mathcal{I}(R)$, the corresponding Green’s class is the entire semigroup:
	\[
	\mathcal{L}_{A}
	= \mathcal{R}_{A}
	= \mathcal{J}_{A}
	= \mathcal{H}_{A}
	= \mathcal{D}_{A}
	= \mathcal{I}(R).
	\]
Therefore, all Green’s relations on $\big(\mathcal{I}(R), \cdot , \preccurlyeq\big)$ are universal.
\end{proof}
Next we provide certain regularity conditions on $\big(\mathcal{I}(R), \cdot , \preccurlyeq\big)$.
\begin{theorem}
	Let $R$ be a commutative principal ideal ring with unity.  
	Then the ordered semigroup
	$\big(\mathcal{I}(R), \cdot , \preccurlyeq \big)$ satisfies the following properties:
	\begin{enumerate}
		\item it is intra-regular
		\item it is an inverse ordered semigroup
		\item every element is an ordered idempotent
		\item it is completely regular
		\item it is group-like
		\item it is a Clifford ordered semigroup.
	\end{enumerate}
\end{theorem}
 
\begin{proof}
	\begin{enumerate}
		\item Consider $a, x, y \in R$, we have $xa^{2}y = a(xay)$, hence $a \mid xa^{2}y$. Therefore, $A \preccurlyeq XA^{2}Y,$
	where $A = \langle a \rangle$, $X = \langle x \rangle$, and $Y = \langle y \rangle$.
	Thus $\big(\mathcal{I}(R), \cdot , \preccurlyeq\big)$ is intra-regular.

	\item For any $a,a' \in R$, we have $a \mid a a' a$ and $a' \mid a' a a'$.
	Consequently, $A \preccurlyeq A A'A
	\quad \text{and} \quad
	A'\preccurlyeq
	A'AA'$, where $A = \langle a \rangle, A' = \langle a' \rangle$. Thus every element is an ordered inverse of every other element and
	$\mathcal{I}(R)$ consists of a single $\mathcal{H}$-class.
	Hence it is an inverse ordered semigroup.
	
	\item Since $a \mid a^{2}$ for all $a \in R$, we obtain $A \preccurlyeq  A^{2}$, where $A = \langle a \rangle$, therefore every element of $\mathcal{I}(R)$ is an ordered idempotent.
	
	\item For all $x \in R$, we have $a \mid a^{2} x a^{2}$ which yields $A \preccurlyeq A^{2}X A^{2}$, where $A = \langle a \rangle, X = \langle x \rangle $.
	Hence $\big(\mathcal{I}(R), \cdot , \preccurlyeq\big)$ is completely regular.
	
	\item For $a,b \in R$, since $a \mid ab$ and $a \mid ba$, we have $A \preccurlyeq AB$ and $A \preccurlyeq BA$, where $A = \langle a \rangle $ and $B = \langle b \rangle$,
	showing that the ordered semigroup is group-like.
	
	\item Since every element of $\mathcal{I}(R)$ is an ordered idempotent
	$ E = \langle e \rangle \in \mathcal{I}(R)$.
	For any $a \in R$, if $ae \mid eua$ and $ea \mid ave$ for all $u,v \in R$, then $AE \preccurlyeq EUA $ and $EA \preccurlyeq AVE $, where $A = \langle a \rangle$, $ U = \langle u \rangle$ and $V = \langle v \rangle$.
\end{enumerate}
	Therefore
	$\big(\mathcal{I}(R), \cdot , \preccurlyeq\big)$ is a Clifford ordered semigroup.
\end{proof}


\begin{example} Ordered semigroup of ideals  $(\mathcal{I}(\mathbb{Z}), \cdot , \preccurlyeq)$; of $\mathbb{Z}$.\\
	Clearly $n$ and $-n$ generate the same principal ideal in $\mathbb{Z}$, thus 
	\[
	\mathcal{I}(\mathbb{Z}) = \{\, \langle n \rangle \mid n \in \mathbb{N} \cup \{0\} \,\}
	\]
	is a semigroup under ideal multiplication.
	In particular for any $\langle m \rangle, \langle n \rangle \in \mathcal{I}(\mathbb{Z})$,
	their product is given by
$	\langle m \rangle \cdot \langle n \rangle = \langle mn \rangle$
	and the order relation $\preccurlyeq$ on $\mathcal{I}(\mathbb{Z})$ by
	\[
	\langle n \rangle \preccurlyeq \langle m \rangle
	\quad \text{if and only if} \quad n \mid m \; \text{in}\, \mathbb{Z}.
	\]
	Let $\mathbb{N}^{0} = \mathbb{N} \cup \{0\}$ and define a mapping
	\[
	\varphi : \mathcal{I}(\mathbb{Z}) \longrightarrow \mathbb{N}^{0},
	\quad 	\varphi(\langle n \rangle) = n.
	\]
	It is straightforward to verify that $\varphi$ is a bijection preserving
	both the semigroup operation and the order relation.
	Hence
	\[
	(\mathcal{I}(\mathbb{Z}), \cdot , \preccurlyeq)
	\cong
	(\mathbb{N}^{0}, \cdot , \mid).
	\]
\end{example}
\begin{example} Ordered semigroup of ideals $(\mathcal{I}(\mathbb{Z}_{n}), \cdot , \preccurlyeq)$ of $\mathbb{Z}_{n}$.\\
	Each  ideal of the ring $\mathbb{Z}_n$ of integers modulo $n$  is generated by a divisor of $n$ (see\cite{punin}). Consequently
	\[
	\mathcal{I}(\mathbb{Z}_n)
	=
	\{\, \langle d \rangle \mid d \text{ is a  divisor of } n \,\}.
	\]
	For any $a \in \mathbb{Z}_n$, $	\langle a \rangle = \langle \gcd(a,n) \rangle $ and for 
	$\langle d_1 \rangle, \langle d_2 \rangle \in \mathcal{I}(\mathbb{Z}_n)$, their product under ideal multiplication is given by
	\[
	\langle d_1 \rangle \cdot \langle d_2 \rangle
	= \langle d_1 d_2 \rangle
	= \langle \gcd(d_1 d_2, n) \rangle.
	\]
	Define a partial order $\preccurlyeq$ on $\mathcal{I}(\mathbb{Z}_n)$ by
	\[
	\langle d_1 \rangle \preccurlyeq \langle d_2 \rangle
	\quad \text{if and only if} \quad
	d_1 \mid d_2 \; \text{ in}\; \mathbb{Z}_n,
	\]
	equivalently $\langle d_1 \rangle \preccurlyeq \langle d_2 \rangle$ if and only if
	$d_2 = m d_1$ for some integer $m$. Since both $d_1$ and $d_2$ divide $n$, the
	integer $m$ is also a divisor of $n$.
	Let $D(n)$ denote the set of all positive divisors of $n$.
	Define a binary operation $*$ on $D(n)$ by
	\[
	d_1 * d_2 = \gcd(d_1 d_2, n),
	\quad \text{for all } d_1, d_2 \in D(n).
	\]
	The divisibility relation defines a partial order on $D(n)$ and this order is compatible with the operation $*$. Thus $(D(n), *, \mid)$ is an ordered semigroup.
	
	The a map
	\[
	\varphi : \mathcal{I}(\mathbb{Z}_n) \longrightarrow D(n),
	\qquad
	\varphi(\langle d \rangle) = d.
	\]
	is well defined, bijective and preserves both the semigroup operation
	and the order relation.
	Hence,
	\[
	(\mathcal{I}(\mathbb{Z}_n), \cdot , \preccurlyeq)
	\cong
	(D(n), *, \mid).
	\]

\end{example}

\begin{example}  Ordered semigroup $(\mathcal{I}(F[x]), \cdot , \preccurlyeq)$ of ideals of $F[x]$.\\

Consider the polynomial ring $F[x]$, over a field $F$. Then
\[
\mathcal{I}(F[x])
=
\{\, \langle f(x) \rangle \mid f(x) \in F[x] \,\}.
\]
is an ideal where the 
ideal multiplication is given by
$\langle f(x) \rangle \cdot \langle g(x) \rangle
= \langle f(x) g(x) \rangle $, for any $\langle f(x) \rangle, \langle g(x) \rangle \in \mathcal{I}(F[x])$
and the order relation $\preccurlyeq$ on $\mathcal{I}(F[x])$ by
\[
\langle f(x) \rangle \preccurlyeq \langle g(x) \rangle
\quad \text{if and only if} \quad
f(x) \mid g(x)\; \text{in}\; F[x]. 
\]
Observe that for any  $ a \in F^{*}  =F \setminus\{0\} $, then 
$\langle a \rangle = F[x]$, while $\langle 0 \rangle = \{0\}$ is the zero
ideal. Moreover two principal ideals $\langle f(x) \rangle$ and
$\langle g(x) \rangle$ coincide if and only if $f(x)$ and $g(x)$ differ by a
nonzero scalar multiple. 
The relation $\sim$ on  $F[x]$, defined by
\[
p(x) \sim q(x)
\iff
p(x) = a q(x) \quad \text{for some } a \in F^{*}.
\]
is an equivalence relation on $F[x]$. and $F[x]/\sim
= \{\, [f(x)] \mid f(x) \}$
denote the set of equivalence classes of $\sim$.
For $[f(x)], [g(x)] \in F[x]/\sim$, define multiplication by
$[f(x)] [g(x)] = [f(x) g(x)].$
Then $F[x]/\sim$ is a semigroup and we introduce a partial order on $F[x]/\sim$ by setting
$[f(x)] \mid [g(x)]
\iff
f(x) \mid g(x) \text{ in }F[x].$
Now the map
\[
\psi:\mathcal{I}(F[x])\longrightarrow F[x]/\sim
\quad \text{by} \quad
\psi(\langle f(x)\rangle)=[f(x)].
\]
is well defined, bijective and preserves both multiplication and
order. Consequently
\[
(\mathcal{I}(F[x]),\cdot,\preccurlyeq)
\cong
(F[x]/\sim,\cdot,\mid).
\]
\end{example}
\section{Category of ideals of a principal ideal ring}\label{pir}

In \cite{calcutta}, we introduced the categories $\mathbb{L}_R$ and $\mathbb{R}_R$, 
whose objects are the left and right ideals respectively of a Noetherian ring $R$ with unity and morphisms left and right $R$-linear maps and, it is  shown that these are 
preadditive categories with a zero object which can be viewed as full subcategories of the category of $R$-modules. Further, these are caategories with subobjects and their morphisms admit factorization. 

In the following, the ring $R$  we discuss is a commutative principal ideal ring with unity, 
so every ideal of $R$ is generated by a single element, hence $R$ is Noetherian. Since $R$ is commutative the notions of left and right ideals coincide, therefore, the categories $\mathbb{L}_R$ and $\mathbb{R}_R$ are identical and we denote this common category by
$\mathbb{I}_R$.  
Each objects in $\mathbb{I}_R$ has the form $A = \langle a \rangle$ for some $a \in R$ and any element in $A$ is written as $ra$ for some $r \in R$. For ideals $A = \langle a \rangle$ and $B = \langle b \rangle$, the set of morphisms $\mathbb{I}_R(A,B)$ consist of all $R$-linear maps from $A$ to $B$, thus, a function $f : A \to B$ is a morphism if and only if
\[
f(r_1 a + r_2 a) = (r_1 + r_2) f(a)
\quad \text{for all } r_1, r_2 \in R,
\]
since each ideal is generated by a single element any such morphism is completely determined by the image of the generator.  
Also $f : A \to B$ is $R$-linear, and $f(a) \in B = \langle b \rangle$ implies there exists $c \in R$ such that $f(a) = cb$, hence, for every $ra \in A,\quad f(ra) = rcb.$
Thus, every morphism from $A$ to $B$ is uniquely determined by a choice of $c \in R$, 
further, being a Noetherian ring, the category $\mathbb{L}_R$ admits subobjects and the subobject relation in $\mathbb{L}_R$ given by the inclusion of left ideals. 

Next we describe the category $\mathbb{I}_R$,  when $R$ is a commutative principal ideal ring and proceed to show that it is a category with subobjects. For $A = \langle a \rangle,\, B = \langle b \rangle \in \mathbb{I}_R$,  define
\[
A \subseteq B \iff a = rb \text{ for some } r \in R
\]
which is clearly a partial order on objects of $\mathbb{I}_R$.
For $A \subseteq B$, the map $j_{(A,B)} : A \to B$ given by  
\[
j_{(A,B)}(ra) = ra, 
\] 
is $R$-linear and hence a morphism in $\mathbb{I}_R$. Moreover, it is a monomorphism and termed it as the inclusion of $A$ into $B$.
Consider the subcategory $\mathcal{P}$ of $\mathbb{I}_R$ having the same objects as $\mathbb{I}_R$ and whose morphisms are precisely these inclusion maps $j_{(A,B)}$, since there is at most one such morphism between any two objects, $\mathcal{P}$ is a strict preorder.
Finally, suppose that $j_{(A,D)}$ and $j_{(B,D)}$ are morphisms in $\mathcal{P}$ and that
\[
j_{(A,D)} = h \, j_{(B,D)}
\]
for some morphism $h : A \to B$ in $\mathbb{I}_R$, as both $j_{(A,D)}$ and $j_{(B,D)}$ are inclusion maps into $D = \langle d \rangle$, it follows that $h$ must coincide with the inclusion map from $A$ into $B$. 
Therefore, the pair $(\mathbb{I}_R, \mathcal{P})$ satisfies the axioms of a category with subobjects.

\section{Correspondence of ideal categories and ordered semigroup structures of  PIR}

In the following we explore the connection between categorical  and ordered semigroup structures arising from the ideals of a commutative principal ideal ring.

Consider the commutative principal ideal ring $R$ with unity. It is already seen that $\mathbb{I}_R$, the category of ideals of $R$ is a category with subobjects and  $\mathcal{I}(R)$ equipped with ideal multiplication and order relation $\preccurlyeq$ defined by divisibility of generators is an ordered semigroup.

To formalize the relation between these structures,  we construct a category $\mathcal{C}$ 
whose objects are those of $\mathcal{I}(R)$, and for $A = \langle a \rangle,\, B = \langle b \rangle \in \mathcal{I}(R) $, define a morphism
\[
A \to B
\quad \text{if and only if} \quad
A \preccurlyeq B,
\]
i.e., whenever $a \mid b$ denoted by $f_{(A,B)}$. Since the divisibility relation determines the existence of a morphism there can be at most one morphism between any two objects.
The composition in $\mathcal{C}$ is induced by the transitivity of divisibility, that is when 
$f_{(A,B)} : A \to B
\quad \text{and} \quad
f_{(B,C)} : B \to C,$
then $a \mid b$ and $b \mid c$, which implies $a \mid c$, accordingly, 
$
f_{(A,B)}\, f_{(B,C)} = f_{(A,C)}.$
 For each object $A = \langle a \rangle$, the relation $a \mid a$ yields the identity morphism $f_{(A,A)}$ . The identity laws follow immediately from the definition of composition and associativity of composition is a direct consequence of the transitivity of divisibility. Thus, $\mathcal{C}$ is a category in which there is at most one morphism between any two objects; ie., it is a strict preorder category.

Now, consider the map
$F : \mathcal{C} \to \mathbb{I}_R$, whose object mad is identity and for a morphism 
$f_{(A,B)} : A \to B$ in $\mathcal{C}$, define
$F(f_{(A,B)}) = j_{(B,A)},$
where $j_{(B,A)} : B \to A$ is the inclusion map in $\mathbb{I}_R$. It can be seen that $F$ is a contravariant functor. For, consider any object $A$, we have
$F(f_{(A,A)}) = j_{(A,A)},$
which is the identity morphism on $A$ in $\mathbb{I}_R$.
Next, consider morphisms $f_{(A,B)} : A \to B,\, f_{(B,C)} : B \to C$, then
$F(f_{(A,B)}) = j_{(B,A)}, \quad F(f_{(B,C)}) = j_{(C,B)}$, and their composition in $\mathbb{I}_R$ is
$j_{(C,B)} \, j_{(B,A)} = j_{(C,A)},$
which is the inclusion of $C$ into $A$. On the other hand,
$F(f_{(A,B)} \circ f_{(B,C)}) = F(f_{(A,C)}) = j_{(C,A)}$, thus,
\[
F(f_{(A,B)} \, f_{(B,C)}) = F(f_{(B,C)}) \, F(f_{(A,B)}),
\]
showing that composition is reversed.

Therefore, we conclude that there is a natural correspondence between the ordered semigroup structure $\mathcal{I}(R)$ and the categorical structure $\mathbb{I}_R$ given by the contravariant functor $F$ from $\mathcal{C}$ to $\mathbb{I}_R$.

	\begin{thebibliography}{99}
	 \bibitem{artin} M.~Artin, \textit{Algebra}, Prentice Hall, New Jersey, ISBN 0-13-004763-5 (1991).
	
	\bibitem{buniya} A.~K.~Bhuniya, K.~Hansda, \textit{On completely regular and Clifford ordered semigroups}, Afrika Mathematica 31, 1029--1045 (2020).
	
	  \bibitem{txtbk} T.~S.~Blyth, \textit{Lattices and Ordered Algebraic Structures}, Springer, London, ISBN 978-1-84628-127-3 (2005).

	\bibitem{cao} Y.~Cao, \textit{Characterization of Regular Ordered Semigroups by Quasi-ideals}, Vietnam Journal of Mathematics 30(3), 239--250 (2002).
	
	
	\bibitem{jamadar} K.~Hansda, A.~Jamadar, \textit{Characterization of inverse ordered semigroups by their ordered idempotents and bi-ideals}, Quasigroups and Related Systems 28, 77--88 (2020).
	
	
	\bibitem{keharegularity} N. Kehayopulu, \textit{On regular duo ordered semigroups}, Math. Japon., 535–540, 37(3) (1992).
	
	\bibitem{intrregular} N. Kehayopulu, \textit{ On intra-regular ordered semigroups} Semigroup Forum 46, 271–278 (1993).
	
	\bibitem{kahaloideal} N.~Kehayopulu, \textit{Ideals and Green's relations in ordered semigroups}, International Journal of Mathematics and Mathematical Sciences 6, 1--8 (2006)
	
	\bibitem{maclane} S.~Mac~Lane, \textit{Categories for a Working Mathematician}, 2nd ed., Springer-Verlag, New York, ISBN 0-387-98403-8 (1998).
	
\bibitem{mpk} P.~K.~Minnumol, P.~G.~Romeo, \textit {On bi-ideals of ordered full transformation semigroups}, Discussiones Mathematicae - General Algebra and Applications 45 (2025) 229–239.
	
	\bibitem{mpkinverse} P.~K.~Minnumol, P.~G.~Romeo, \textit{Inverse Transversal of Ordered Semigroups}, Asian-European Journal of Mathematics (2025).

	
	\bibitem{musili} C.~Musili, \textit{Introduction to Rings and Modules}, Narosa Publishing House, New Delhi, ISBN 81-85198-64-0 (1992).
	
	\bibitem{kss} K.~S.~S.~Nambooripad, \textit{Theory of Cross Connections}, Publication No.~28, Centre for Mathematical Sciences, Trivandrum (1994).
	
	\bibitem{punin} W.~Puninagool, J. Sanwong \textit{Ideals of the Multiplicative Semigroups Zn and their Products}, Kyungpook
	Mathematical Journal, 41-46, 49(2009).
	
	\bibitem{calcutta} P.~G.~Romeo, P.~K.~Minnumol, \textit{Ideal Categories of Noetherian Rings}, Bulletin of the Calcutta Mathematical Society 116(4), 467--476 (2024).
	

\end{thebibliography}
\end{document}





=======================                   ===================================


existence of identity morphisms follows from the fact that every element divides itself.


Thus, the inclusion $A \subseteq B$ in the category $\mathbb{I}_R$ corresponds to the reversed order $B \preccurlyeq A$ in the ordered semigroup $\mathcal{I}(R)$. This reversal highlights the contravariant nature of the correspondence between these two structures.


Let $A = \langle a \rangle$ and $B = \langle b \rangle$ be ideals of $R$. Then $A \subseteq B$ if and only if there exists $r \in R$ such that $a = rb$. Equivalently, this condition can be written as $b \mid a$. Hence, inclusion of ideals is completely characterized by divisibility among their generators.

Now consider the set  of all ideals of $R$ . As established earlier this gives $\mathcal{I}(R)$ the structure of . In this setting, the condition $b \mid a$ is equivalent to
$ B \preccurlyeq A.



\begin{theorem}(cf.\cite{calcutta})\label{theorensuboject}
	The category $ \mathbb{L}_{R} $,  of all left ideals of a Noetherian ring $ R $ 
	is a category with subobjects and every morphisms have canonical factorization.
\end{theorem}


\subsection{Category of ideals of a principal ideal ring}\label{pir}

In \cite{calcutta}, the categories $\mathbb{L}_R [\mathbb{R}_R]$ are described, whose objects are the left [right] ideals of a Noetherian ring $R$ with unity, and whose morphisms are left [right]  $R$-linear transformations. It is further shown that these are preadditive categories with a zero object and are full subcategories of the category of $R$-modules. Moreover, they admit subobjects, and their morphisms satisfy the factorization property.

In this section, we study the ideal category associated with a special class of Noetherian rings, namely commutative principal ideal rings.

Let $R$ be a commutative principal ideal ring with unity. Since every ideal of $R$ is generated by a single element $R$ is a Noetherian ring. Because $R$ is commutative the notions of left ideals and right ideals coincide. Hence the categories of left ideals and right ideals are identical and we denote this common ideal category by
\[
\mathbb{L}_R=\mathbb{R}_R=\mathbb{I}_R .
\]

The objects of the category $\mathbb{I}_R$ are the ideals of the ring $R$. Since $R$ is a principal ideal ring, every ideal is of the form 
$ A = \langle a\rangle, \, a\in R.$
%Let $I=\langle a\rangle$ and $J=\langle b\rangle$.
Any element of $A = \langle a\rangle$ can be written as $ra$ for some $r\in R$. For objects $A = \langle a\rangle, B = \langle b\rangle\in \nu(\mathbb{I}_R)$ we define the set of morphisms from $A$ to $B$, denoted by $\mathbb{I}_R(A, B)$ to be the collection of all $R$-linear transformations from $A$ to $B$. Thus a mapping $f:A\to B$ is a morphism in $\mathbb{I}_R$ if and only if
\[
f(r_{1}a+r_{2}a)=[r_{1}+r_{2}]f(a)
\]
for all  $r_{1},r_{2}\in R$.

Because every ideal of a principal ideal ring is generated by a single element such a morphism is completely determined by the image of the generator. If $f:A\to B$ is an $R$-linear transformation then
$f(ra)=r\,f(a) .$
Since $f(a)\in B=\langle b\rangle$ there exists $c\in R$ such that
$f(a)=cb .$
Consequently for every $ra\in A$ we have
$f(ra)=rcb .$
Thus every morphism from $A$ to $B$ arises from the choice of an element $c\in R$ which determines the image of the generator $a$ in the ideal $B$.

We have seen in Theorem~\ref{theorensuboject} that if $R$ is a Noetherian ring, then the category $\mathbb{L}_{R}$ of left ideals of $R$ is a category with subobjects. We now examine the subobject structure of the ideal category $\mathbb{I}_{R}$ in the special case where $R$ is a commutative principal ideal ring with unity. To determine this structure we construct a suitable subcategory $\mathcal{P}$ of $\mathbb{I}_{R}$ that satisfies the axioms of a category with subobjects.

%Let $\nu(\mathbb{I}_R)$ denote the class of objects of $\mathbb{I}_R$.
For any two ideals $A = \langle a\rangle, B = \langle b\rangle\in \nu(\mathbb{I}_R)$, we define a partial order by  whenever the ideal generated by $a$ is contained in the ideal generated by $b$. That is
\begin{center}
	$A \subseteq B \iff a=rb $ for some $r\in R$.
\end{center}
For each pair $ A \subseteq B  $, %define a morphism 
$j_{(A,B)}:A\to B$ denote the the inclusion map defined 
by
\[
j_{(A,B)}(a)= a.% \qquad  r\in R.
\]
Thus the inclusion map sends the element $ra\in A$ to the same element $ra$, now regarded as an element of $B$. This map clearly preserves addition and multiplication by elements of $R$ and hence it is an $R$-linear transformation. Therefore $j_{(A,B)}$ is a morphism in the category $\mathbb{I}_R$. Moreover it is a monomorphism and uniquely represents the inclusion of $A$ into $B$.

Let $\mathcal{P}$ denote the subcategory of $\mathbb{I}_R$ whose objects coincide with those of $\mathbb{I}_R$ and whose morphisms are precisely these inclusion maps $j_{(A,B)}$. Since there is at most one inclusion map between any two ideals the subcategory $\mathcal{P}$ forms a strict preorder.

Finally suppose that $j_{(A,D)}$ and $j_{(B,D)}$ are morphisms in $\mathcal{P}$ and that
\[
j_{(A,D)} = h\, j_{(B,D)}
\]
for some morphism $h: A \to B $ in $\mathbb{I}_R$. Since both $j_{(A,D)}$ and $j_{(B,D)}$ are inclusion maps into $ D = \langle d\rangle$ it follows that the morphism $h$ must also be the inclusion map from $A$ into $B$. Therefore the pair $(\mathbb{I}_R,\mathcal{P})$ satisfies the axioms of a category with subobjects.

\section{Correspondence between Categorical and ordered semigroup structures of ideals in ring}

In this section we examine the connection between the categorical structure and the ordered semigroup structure of ideals in a commutative principal ideal ring.

Let $R$ be a commutative principal ideal ring with unity and  $\mathbb{I}_R$ denote the category of ideals of $R$.
In the above section we proved that  $\mathbb{I}_R$ is a category with subobjects. In the category $\mathbb{I}_R$ subobjects coincide with inclusions of ideals. %Let $\langle a\rangle$ and $\langle b\rangle$ be two objects of $\mathbb{I}_R$. By the characterization of subobjects in this category,\[\langle a\rangle \text{ is a subobject of } \langle b\rangle\iff\langle a\rangle \subseteq \langle b\rangle \]

%Since $R$ is a commutative principal ideal ring every ideal is generated by a single element. Consequently
For ideals $A = \langle a \rangle$ and $B = \langle b \rangle$ the inclusion $A \subseteq B $ holds if and only if there exists an element $r \in R$ such that
$a = rb.$
Equivalently this condition can be expressed as $b \mid a$. Thus in the ideal category $\mathbb{I}_R$ the subobject relation is completely determined by the divisibility relation between generators of ideals.

On the other hand consider the collection $\mathcal{I}(R)$ of all ideals of $R$ equipped with ideal multiplication as the semigroup operation and the divisibility relation as a partial order. As shown earlier in Section~\ref{section} this structure turns $\mathcal{I}(R)$ into an ordered semigroup. From this perspective, the condition $b \mid a$ is precisely the order relation
$B \preccurlyeq A .$

Therefore whenever $A$ is a subobject of $B$ in the ideal category $\mathbb{I}_R$ the corresponding elements in the ordered semigroup $\mathcal{I}(R)$ appear in the reverse order namely $ B \preccurlyeq A$. This reversal reflects the contravariant nature of ideal inclusion when interpreted through divisibility.\\
This phenomenon can also be described in categorical terms. Let $\mathcal{C}$ be a category whose object set is $\nu\mathcal{C}=\mathcal{I}(R)$ the collection of all ideals of $R$. Thus each object of $\mathcal{C}$ is of the form $ A = \langle a\rangle$ for some $a\in R$. For two objects $A = \langle a\rangle$ and $ B = \langle b\rangle$, we define a morphism
\begin{center}
	$A \longrightarrow B$
	if and only if $A \preccurlyeq B $ in $\mathcal{I}(R)$
\end{center} i.e.,$a \mid b$. This morphism is denoted by $f_{(A,B)}$. Since the divisibility relation only determines the existence of a morphism there can be at most one morphism between any two objects.

Composition of morphisms corresponds to the transitivity of divisibility. Let $A = \langle a\rangle, B =\langle b\rangle $  and $ C = \langle c\rangle \in \nu \mathcal{C}$. If
$f_{(A,B)} : A \to B
\quad \text{and} \quad
f_{(B,C)} : B \to C $
are morphisms in $\mathcal{C}$ then $a \mid b$ and $b \mid c$. Hence $a \mid c$, and we define the composite morphism by
\[
f_{(A,B)}\,\,f_{(B,C)} = f_{(A,C)}.
\]

Identity morphisms arise from the fact that every element divides itself. Indeed since $a \mid a$ there exists a morphism $f_{(A,A)} : A \to  A $ which serves as the identity morphism of the object $A$. For  morphisms $f_{(A,B)} :  A \to B $ and $f_{(X,A)} : X \to A $ we have
\[
f_{(X,A)}\,f_{(A,A)} = f_{(X,A)}
\quad \text{and} \quad
f_{(A,A)}\, f_{(A,B)} = f_{(A,B)}.
\]
Thus the identity law holds.

Associativity of composition follows from the transitivity of divisibility. If
$f_{(A,B)} : A \to B$,
$f_{(B,C)} : B \to C$ and
$f_{(C,D)} : C \to D$
are morphisms in $\mathcal{C}$, then 
\[
(f_{(A,B)}\,f_{(B,C)})\,f_{(C,D)}
\, = \,
f_{(A,B)}\,(f_{(B,C)}\,f_{(C,D)}) \, = \, f_{(A,D)}.
\] Hence composition is associative.

Thus the morphisms of $\mathcal{C}$ are determined by the partial order given by the divisibility relation and there is at most one morphism between any two objects. Therefore $\mathcal{C}$ is a strict preorder category.

Now define a mapping
$$F : \mathcal{C} \to \mathbb{I}_R$$
as follows. On objects $F$ acts as the identity. For each morphism 
$f_{(A,B)} \in \mathcal{C}(A,B)$, define
\[
F(f_{(A,B)}) = j_{(B,A)}
\]
where 
$j_{(B,A)} : B \to A$
is the natural inclusion map in $\mathbb{I}_R$.

We now verify that $F$ is a contravariant functor. First consider identity morphisms. For any object $ A = \langle a\rangle$ in $\mathcal{C}$, since $a \mid a$ the identity morphism is given by $f_{(A,A)}$. Hence
\[
F(f_{(A,A)}) = j_{(A,A)}
\]
which is precisely the identity map on $A$ in the category $\mathbb{I}_R$. Thus $F$ preserves identity morphisms.

Next suppose that $A = \langle a\rangle, B =\langle b\rangle , C = \langle c\rangle \in \nu \mathcal{C}$ and   
$f_{(A,B)} : A \to B $ and $
f_{(B,C)} : B \to C$
are morphisms in $\mathcal{C}$. Then $a \mid b$ and $b \mid c$ which implies $a \mid c$ and therefore
$ f_{(A,B)}\,f_{(B,C)} = f_{(A,C)}.$
Applying $F$ we obtain
$F(f_{(A,B)}\,f_{(B,C)}) = F(f_{(A,C)}) = j_{(C,A)}.$
On the other hand
\[
F(f_{(A,B)}) = j_{(B,A)}
\quad \text{and} \quad
F(f_{(B,C)}) = j_{(C,B)}.
\]
Since $j_{(C,B)}$ is the inclusion of $C$ into $B $ and $j_{(B,A)}$ is the inclusion of $B$ into $A$, their composition is the inclusion of $C$ into $A$. Hence
$ j_{(C,B)}\,j_{(B,A)} = j_{(C,A)}$ and $ F(f_{(A,C)}) = F(f_{(B,C)})\, F(f_{(A,B)}) $.
Thus
\[
F(f_{(A,B)}\,f_{(B,C)}) = F(f_{(B,C)})\, F(f_{(A,B)}).
\]
showing that the direction of composition is reversed.
Therefore
$F : \mathcal{C} \to \mathbb{I}_R$
is a contravariant functor. This establishes a natural connection between the categorical structure of ideals in $\mathbb{I}_R$ and the ordered semigroup structure of ideals in $\mathcal{I}(R)$ thereby providing a bridge between category-theoretic and semigroup-theoretic approaches to the study of ideals in commutative principal ideal rings.




%Now we show compatibility with multiplication.
Let $A, B, C \in \mathcal{I}(R)$ with 
$A = \langle a \rangle$, $B = \langle b \rangle$ and $C = \langle c \rangle$, and suppose $A \preccurlyeq B$. 
Then $a \mid b$, so $b = ax$ for some $x \in R$. Hence $
bc = (ax)c = x(ac)$ 
which shows that $ac \mid bc$.

Since 
$AC = \langle a \rangle \langle c \rangle = \langle ac \rangle, \quad
BC = \langle b \rangle \langle c \rangle = \langle bc \rangle.$
Thus $AC \preccurlyeq BC$. Similarly, $CA \preccurlyeq CB$. 





\subsection{Examples}
We next consider examples of ideal semigroups associated with principal
ideal rings.










